% ============================================================================================
% PW-1 item 9 -- Shapes3D rendered success/failure figures (pw1-render-shapes).
% NOT INSTALLED. A later spoke places these in Sec.~\ref{sec:shapes-t7}, after Tab.~\ref{tab:t7}.
% They replace the w19 slots fig:t7-per-arm, fig:t7-rival-wins, fig:t7-rival-wins-s1
% (removed by w22 draft-rulings Ruling 2; see the comment block at DRAFT-v2.tex:477-480).
% Figure files: outputs/pw1-render-shapes/build/a/rc_pw1_shapes_per_model.pdf
%               outputs/pw1-render-shapes/build/a/rc_pw1_shapes_baseline_wins.pdf
% Every number below: outputs/pw1-render-shapes/REQUESTS-pw1.json (key named per line).
% Requires amssymb (\checkmark); the drafts load it.
% ============================================================================================

% ---- (1) the paragraph that introduces the two figures (Shapes results section) ------------
% SRCBLOCK: outputs/pw1-render-shapes/REQUESTS-pw1.json; outputs/pw1-render-shapes/factor-miss-seed0.txt
\paragraph{What the counts look like.} Figs.~\ref{fig:shapes-per-model} and~\ref{fig:shapes-baseline-wins}
show single held-out requests from Tab.~\ref{tab:t7} as images, at seed 0 and $\Kstar=32$, against the
four data-matched baselines of that table (MLP and continuous MLP at $10\times$ training steps, LSTM and
Transformer at their selected schedule, all fitted on the same 131,072 rows) and the lookup table as the
floor. Every model's output on a request is drawn through the same frozen decoder, and each panel
carries the judge's verdict on that image. Over all 5,688 held-out requests at seed 0 the judge marks
the n-CLU valid on 5,177, the MLP on 4,361, the LSTM on 3,703, the Transformer on 2,786, the continuous
MLP on 3,110 and the lookup table on 897.
% SRC: REQUESTS-pw1.json counts_seed0.per_column.<column>.valid; n_heldout_requests
Where the n-CLU is invalid, the image is almost always wrong in one factor only (497 of its 511 invalid
renders), most often the scale (329) or the orientation (125), less often the shape (71).
% SRC: factor-miss-seed0.txt (from scan/verdicts-seed0.json)
Requests on which a baseline succeeds and the n-CLU fails exist for every baseline (310 requests at
seed 0 have at least one); they are not confined to the capsule: 170 of the 310 are capsules, 100
cylinders and 40 spheres.
% SRC: REQUESTS-pw1.json counts_seed0.any_baseline_valid_and_n_clu_invalid
The rows of both figures are drawn by a fixed-seed rule, stated in the captions, never by appearance.
% ENDSRC

% ---- (2) Figure: per-model success and failure grid ---------------------------------------
% SOURCE: outputs/pw1-render-shapes/build/a/rc_pw1_shapes_per_model.pdf (sha256 in SHA256-build1.txt)
\begin{figure}[t]\centering
\includegraphics[width=0.9\linewidth]{figures/rendered/rc_pw1_shapes_per_model.pdf}
\caption{\textbf{Each model succeeding and failing on held-out combinations (Shapes3D, seed 0,
$\Kstar=32$).} Columns: the n-CLU, the four data-matched baselines of Tab.~\ref{tab:t7} (MLP and
continuous MLP at $10\times$ training steps, LSTM and Transformer at their selected schedule) and the
lookup table as the floor. For each model (bar at left) one held-out request it renders validly
(success) and one it renders invalidly (failure); the framed panel is the model the row was drawn for.
Every model's image on every request comes from the same frozen decoder, and the mark is the judge's
verdict on that image ($\checkmark$ valid, $\times$ invalid). The requested values are printed left of
each row (hues as fractions of the hue circle, scale, orientation in degrees); a request either holds out
the object hue $\times$ shape pair or all six factors jointly. \emph{Selection:} a fixed-seed draw
(seed 20260925) among all qualifying held-out requests, under a diversity rule: the target shape is
cycled over the rows so the capsule is never drawn more often than any other shape, and among requests
of the target shape the draw prefers object, floor and wall hues not yet shown. The judge's six read-back
values for every panel are listed in the supplementary material.}
% SRC: REQUESTS-pw1.json select_seed; per_model.rows[*]; carried_fit_seed0; selection_rule
\label{fig:shapes-per-model}
\end{figure}

% ---- (3) Figure: baseline succeeds, n-CLU fails --------------------------------------------
% SOURCE: outputs/pw1-render-shapes/build/a/rc_pw1_shapes_baseline_wins.pdf (sha256 in SHA256-build1.txt)
\begin{figure}[t]\centering
\includegraphics[width=\linewidth]{figures/rendered/rc_pw1_shapes_baseline_wins.pdf}
\caption{\textbf{Held-out requests on which a baseline succeeds and the n-CLU fails (seed 0).} Such
requests exist for every baseline: of the 5,688 held-out requests, the MLP is valid and the n-CLU
invalid on 200, the LSTM on 197, the Transformer on 110 and the continuous MLP on 167 (310 requests have
at least one such baseline; at seeds 1 and 2 there are 332 and 239). They are not all capsules: of the
310 at seed 0, 170 are capsules, 100 cylinders and 40 spheres. One request per baseline is shown; the
n-CLU's panel and the baseline's panel are framed. Same frozen decoder, judge marks, layout and
fixed-seed diversity rule (seed 20260925) as Fig.~\ref{fig:shapes-per-model}.}
% SRC: REQUESTS-pw1.json counts_seed0.baseline_valid_and_n_clu_invalid.<baseline>.n;
%      counts_seed0.any_baseline_valid_and_n_clu_invalid; counts_seeds_1_2_not_drawn.{1,2}.any_baseline_valid_and_n_clu_invalid.n
\label{fig:shapes-baseline-wins}
\end{figure}
